\subsection{正则代数中基域的扩张（特征非零）}

设 \(k\) 为环，\(\rho: A \to B\) 为 \(k\)-代数同态。同态 \(\rho\) 诱导出 A-线性映射 \(\Omega(\rho): \Omega_k(A) \to \Omega_k(B)\)，从而诱导出 B-线性映射 \(\Omega_0(\rho): B \otimes_A \Omega_k(A) \to \Omega_k(B)\) (A, III, 第 135 页)。设 \(\mathbf{T} = (\mathrm{T}_i)_{i \in I}\) 为不定元族，\(\mathbf{t} = (t_i)_{i \in I}\) 为 B 的元素族；对每个多项式 \(f = \sum_{\alpha \in \mathbf{N}^{(I)}} c_\alpha \mathbf{T}^\alpha\)（A{[} \(\mathbf{T}\) {]} 中的），以 \(d^{A} f(\mathbf{t})\) 表示元素 \(\sum_{\alpha} \mathbf{t}^\alpha \otimes dc_\alpha\)（\(B \otimes_A \Omega_k(A)\) 中的）。

引理 1.—设 A-代数 B 有一个生成元族 \(\mathbf{t} = (t_i)_{i\in I}\)，它们满足关系 \(f_\lambda \in \mathrm{A}[\mathbf{T}]\)（\(\lambda \in \Lambda\)）。则 B-线性同态

\[
\psi : (\mathrm{B} \otimes_ {\mathrm{A}} \Omega_ {k} (\mathrm{A})) \oplus \mathrm{B} ^ {(\mathrm{I})} \longrightarrow \Omega_ {k} (\mathrm{B})
\]

由 \(\psi(\alpha,(b_i)) = \Omega_0(\rho)(\alpha) + \sum_{i\in I}b_i\,dt_i\) 定义的映射是满射；其核由诸元素 \(n_{\lambda} = \left(d^{A}f_{\lambda}(\mathbf{t}),\left(\frac{\partial f_{\lambda}}{\partial\mathrm{T}_{i}}(\mathbf{t})\right)_{i\in I}\right)\)（\(\lambda \in \Lambda\)）生成。

考察 B-模与 B-线性映射的序列

\[
\mathrm{B} ^ {(\Lambda)} \xrightarrow {\varphi} (\mathrm{B} \otimes_ {\mathrm{A}} \Omega_ {k} (\mathrm{A})) \oplus \mathrm{B} ^ {(\mathrm{I})} \xrightarrow {\psi} \Omega_ {k} (\mathrm{B}) \longrightarrow 0,
\]

其中 \(\varphi\) 是满足 \(\varphi(e_{\lambda}) = n_{\lambda}\) 的同态；尚需证明该序列是正合的。由 A, II, 第 36 页, 定理 1，只需证明：对每个 B-模 M，序列

\[
0 \to \operatorname{Hom} _ {\mathrm{B}} (\Omega _{k}(\mathrm{B}), \mathrm{M}) \xrightarrow {\operatorname{Hom}(\psi,1)} \operatorname{Hom} _{\mathrm{B}} ((\mathrm{B} \otimes_{\mathrm{A}} \Omega _{k}(\mathrm{A})) \oplus \mathrm{B}^{(\mathrm{I})}, \mathrm{M}) \xrightarrow {\operatorname{Hom}(\varphi,1)} \operatorname{Hom}_{\mathrm{B}}(\mathrm{B}^{(\Lambda)}, \mathrm{M})
\]

是正合的。鉴于微分模的泛性质 (A, III, 第 134 页)，该序列与

\[
0 \to \mathrm{D} _ {k} (\mathrm{B}, \mathrm{M}) \xrightarrow {\psi^ {\prime}} \mathrm{D} _ {k} (\mathrm{A}, \mathrm{M}) \oplus \mathrm{M} ^ {\mathrm{I}} \xrightarrow {\varphi^ {\prime}} \mathrm{M} ^ {\Lambda}
\]

恒同，其中 \(\psi^{\prime}(\mathrm{D}) = (\mathrm{D}\circ \rho,(\mathrm{D}(t_{i}))_{i\in I})\)，\(\varphi^{\prime}(\Delta,(m_{i})) = \big(f_{\lambda}^{\Delta}(\mathbf{t}) + \sum_{i}\frac{\partial f_{\lambda}}{\partial\mathrm{T}_{i}}(\mathbf{t})m_{i}\big)_{\lambda \in \Lambda}\)（依照 A, V, 第 121 页，对每个多项式 \(f = \sum_{\alpha \in \mathbf{N}^{(I)}}c_{\alpha}\mathbf{T}^{\alpha}\)（A{[}T{]} 中的），以 \(f^{\Delta}(\mathbf{t})\) 表示元素 \(\sum_{\alpha}\mathbf{t}^{\alpha}\Delta(c_{\alpha})\)）。该序列的正合性由同上引文命题 1 得出，并注意到：导子 D : B → M 是 \(k\)-线性的，必须且只需 \(D\circ\rho\) 如此。

设 A 为环。A 上存在唯一的 \(\mathbf{Z}\)-代数结构；我们以 \(\Omega(A)\) 简记 A-模 \(\Omega_{\mathbf{Z}}(\mathrm{A})\)。若 \(\rho : k \to \mathrm{A}\) 是环同态，则有 A-模的典范正合序列 (A, III, 第 136 页, 命题 21) \(\mathrm{A} \otimes_k \Omega(k) \to \Omega(\mathrm{A}) \to \Omega_k(\mathrm{A}) \to 0\)。

设 A 包含一个子域，并设 P 为 A 的素子域；则 \(\Omega(\mathrm{P})\) 为零，且 A-模的典范同态 \(\Omega(\mathrm{A}) \to \Omega_{\mathrm{P}}(\mathrm{A})\) 是双射。若此外 A 的特征为 \(p \neq 0\)（依定义，这表示 \(p\) 是素数，\(p1_{\mathrm{A}} = 0\)，且 \(1_{\mathrm{A}} \neq 0\)），则 P 与 \(\mathbf{F}_p\) 恒同。此外，A 的每个导子都在子环 \(\mathbf{A}^p\) 上为零；对 A 的每个子环 \(k\)（含于 \(\mathbf{A}^p\) 中的；特别地，对 A 的每个完全域 \(k\)），典范映射 \(\Omega(\mathrm{A}) \to \Omega_k(\mathrm{A})\) 是双射。

设 A 为特征 \(p \neq 0\) 的环，\((f_i)_{1 \leqslant i \leqslant n}\) 为 A 的元素的一个有限序列。以 \(A_n\) 表示多项式环 \(A[T_1, \ldots, T_n]\) 模由多项式 \(T_i^p - f_i\)（\(1 \leqslant i \leqslant n\)）生成的理想所得的商环。

引理 2.—设环 A 是局部的且 Noether 的。则 \(A_{n}\) 是局部的且 Noether 的。下列条件等价：

\begin{enumerate}
\def\labelenumi{(\roman{enumi})}
\item
  \(A_{n}\) 正则；
\item
  A 正则，且元素 \(1 \otimes df_{i}\)（\(\kappa_{A}\)-向量空间 \(\kappa_{A} \otimes_{A} \Omega(A)\) 中的）线性无关。
\end{enumerate}

环 \(A_{n}\) 是 Noether 的 (III, § 2, 定理 2 的推论 3)。A-模 \(A_{n}\) 是自由的，故是忠实平坦的；若 \(A_{n}\) 正则，则 A 正则 (§ 4, 第 5 节, 命题 8 b))。我们对 n 用归纳法推理，引理在 n = 0 时是显然的。

\begin{enumerate}
\def\labelenumi{\Alph{enumi})}
\tightlist
\item
  先处理 n = 1 的情形，令 \(T_{1} = T\)，\(f_{1} = f\)。以 a 表示 f 在 \(\kappa_{A}\) 中的类，并区分两种情形：a 属于或不属于 \(\kappa_{A}^{p}\)。若 \(a \notin \kappa_{A}^{p}\)，则多项式 \(T^{p} - a\) 在 \(\kappa_{A}\) 中不可约 (A, V, 第 24 页, 引理 1)，而 \(\kappa_{A} \otimes_{A} A_{1}\) 同构于域 \(\kappa_{A}[T]/(T^{p} - a)\)。于是理想 \(m_{A}A_{1}\)（\(A_{1}\) 的）是极大的，从而环 \(A_{1}\) 是局部的 (V, § 2, 第 1 节, 命题 1)。若 A 正则，则 \(A_{1}\) 正则 (VIII, § 5, 第 1 节, 命题 1)。据 A, V, 第 99 页, 命题 6，\(\Omega(\kappa_{A})\) 中的元素 da 不为零；由于它是典范映射 \(\kappa_{A} \otimes_{A} \Omega(A) \to \Omega(\kappa_{A})\) 下 \(1 \otimes df\) 的像，后者不为零。这就在此情形证明了引理。
\end{enumerate}

现在设 \(a\) 属于 \(\kappa_{\mathrm{A}}^{p}\)。于是存在 A 的元素 \(g\) 使 \(f - g^{p} \in \mathfrak{m}_{\mathrm{A}}\)。令 \(h = f - g^{p}\)。由于 \(\mathrm{T}^{p} - f = (\mathrm{T} - g)^{p} - h\)，A-代数 \(A_{1}\) 同构于 \(\mathrm{A}[\mathrm{T}] / (\mathrm{T}^{p} - h)\)。据 VIII, § 5, 第 4 节, 命题 4，环 \(A_{1}\) 是局部的，且为使它正则，必须且只需 A 正则且 \(h\) 不属于 \(\mathfrak{m}_{\mathrm{A}}^{2}\)。而由于 \(\kappa_{A}\) 在素域上形式光滑 (第 3 节, 定理 1)，典范映射

\[
\bar {d}: \mathfrak {m} _ {\mathrm{A}} / \mathfrak {m} _ {\mathrm{A}} ^ {2} \to \kappa_ {\mathrm{A}} \otimes_ {\mathrm{A}} \Omega (\mathrm{A})
\]

是单射 (第 2 节, 注 1)；而 \(\bar{d}\) 把 \(h\) 模 \(\mathfrak{m}_{\mathrm{A}}^{2}\) 的类映为 \(1 \otimes dh = 1 \otimes d(f - g^{p}) = 1 \otimes df\)。这就在第二种情形证明了引理，并完成了 \(n = 1\) 情形的证明。

\begin{enumerate}
\def\labelenumi{\Alph{enumi})}
\setcounter{enumi}{1}
\tightlist
\item
  设 \(n > 1\)。环 \(\mathrm{A}_{1}\) 由已处理的情形是局部的且 Noether 的。\(\mathrm{A}_{1}\)-代数 \(\mathrm{A}_{n}\) 与 \(\mathrm{A}_{1}[\mathrm{T}_{2},\ldots ,\mathrm{T}_{n}]\) 模由诸 \(\mathrm{T}_i^p -f_i\)（\(i\geqslant 2\)）生成的理想所得的商恒同；由归纳假设，它是一个局部环，且条件 (i) 等价于下列两个条件的合取：
\end{enumerate}

(i') A \(_{1}\) 正则；

(i'\,') 元素 \(1 \otimes df_2, \ldots, 1 \otimes df_n\)（\(\kappa_{\mathrm{A}_1}\)-向量空间 \(\kappa_{\mathrm{A}_1} \otimes_{\mathrm{A}_1} \Omega(\mathrm{A}_1)\) 中的）线性无关。

而 (i') 由刚才所见等价于

(ii') A 正则，且元素 \(1 \otimes df_{1}\)（\(\kappa_{A}\)-向量空间 \(\kappa_{A} \otimes_{A} \Omega(A)\) 中的）不为零。

由引理 1，典范同态 \(\mathrm{A}_{1} \otimes_{\mathrm{A}} \Omega(\mathrm{A}) \to \Omega(\mathrm{A}_{1})\) 诱导出 \(((\mathrm{A}_{1} \otimes_{\mathrm{A}} \Omega(\mathrm{A})) / \mathrm{A}_{1}(1 \otimes df_{1})) \oplus \mathrm{A}_{1}\) 到 \(\Omega(\mathrm{A}_{1})\) 上的一个同构，从而诱导出 \((\kappa_{\mathrm{A}_{1}} \otimes_{\mathrm{A}} \Omega(\mathrm{A})) / \kappa_{\mathrm{A}_{1}}(1 \otimes df_{1})\) 到 \(\kappa_{\mathrm{A}_{1}} \otimes_{\mathrm{A}_{1}} \Omega(\mathrm{A}_{1})\) 中的一个单同态。由于 \(\kappa_{\mathrm{A}_{1}} \otimes_{\mathrm{A}} \Omega(\mathrm{A})\) 与 \(\kappa_{\mathrm{A}_{1}} \otimes_{\kappa_{\mathrm{A}}} (\kappa_{\mathrm{A}} \otimes_{\mathrm{A}} \Omega(\mathrm{A}))\) 恒同，断言 (i'') 于是等价于：

(ii'') 元素 \(1 \otimes df_2, \ldots, 1 \otimes df_n\) 在 \((\kappa_{\mathrm{A}} \otimes_{\mathrm{A}} \Omega(\mathrm{A})) / \kappa_{\mathrm{A}}(1 \otimes df_1)\) 中线性无关。

而 (ii') 与 (ii'') 的合取等价于 (ii)，这就证明了引理。

命题 6. 设 \(k\) 为特征 \(p \neq 0\) 的域，\(k'\) 为 \(k\) 的一个有限次且高度 \(\leqslant 1\) 的纯不可分扩张，A 为局部正则 \(k\)-代数。则 \(\mathrm{A}_{(k')}\) 是局部环，且下列条件等价：

\begin{enumerate}
\def\labelenumi{(\roman{enumi})}
\item
  环 \(A_{(k')}\) 正则；
\item
  \(\kappa_{A}\)-线性映射
\end{enumerate}

\[
\kappa_ {\mathrm{A}} \otimes_ {k ^ {\prime p}} \Omega_ {k ^ {p}} (k ^ {\prime p}) \longrightarrow \kappa_ {\mathrm{A}} \otimes_ {\mathrm{A}} \Omega (\mathrm{A})
\]

由典范单射 \(k^{\prime p}\to \mathrm{A}\) 导出的映射是单射。

实际上，设 \((x_{i})_{i\in \mathrm{I}}\) 为一个有限 \(p\)-基（\(k'\) 在 \(k\) 上的） (A, V, 第 98 页)；对每个 \(i\in \mathrm{I}\)，令 \(f_{i} = x_{i}^{p}\in k\)。\(k\)-代数 \(k'\) 与 \(k[(\mathrm{T}_i)_{i\in \mathrm{I}}]\) 模由多项式 \(\mathrm{T}_i^p - f_i\) 生成的理想所得的商恒同，因而 A-代数 \(\mathrm{A}_{(k')}\) 是 \(\mathrm{A}[(\mathrm{T}_i)_{i\in \mathrm{I}}]\) 模由多项式 \(\mathrm{T}_i^p - f_i1_{\mathrm{A}}\) 生成的理想所得的商。

此外，\((f_{i})_{i\in\mathrm{I}}\) 是 \(k^{\prime p}\) 在 \(k^{p}\) 上的一个 p-基，且 \(k^{\prime p}\)-向量空间 \(\Omega_{k^{p}}(k^{\prime p})\) 以诸 \(df_{i}\) 组成的族为一个基 (A, V, 第 97 页, 定理 1)。命题 6 于是由引理 2 得出。

